\documentclass{article}
\usepackage{amsmath,amssymb,amsthm}
\usepackage{graphicx}
\usepackage{hyperref}
\usepackage{algorithm2e}
\usepackage{booktabs}

\theoremstyle{definition}
\newtheorem{theorem}{Theorem}
\newtheorem{lemma}[theorem]{Lemma}
\newtheorem{proposition}[theorem]{Proposition}
\newtheorem{corollary}[theorem]{Corollary}
\newtheorem{definition}[theorem]{Definition}
\newtheorem{remark}[theorem]{Remark}
\newtheorem{example}[theorem]{Example}

\DeclareMathOperator{\rank}{rank}

\title{Hadamard Factorization of $4 \times 4$ Matrices: \\ From Finite Fields to the Reals}
\author{Igor Rivin}
\date{\today}

\begin{document}

\maketitle

\begin{abstract}
We investigate when a $4 \times 4$ matrix $C$ can be expressed as the Hadamard (element-wise) product of two rank-2 matrices: $C = A \circ B$ with $\rank(A), \rank(B) \leq 2$. Through exhaustive computational search, we characterize the obstruction in $\mathbb{F}_2$ (26.3\% of full-rank matrices are not expressible), extend to $\mathbb{F}_3$ (where the pattern reverses dramatically), and develop efficient algorithms for the real case. We discover a surprising ``U-shaped'' factorizability curve over $\mathbb{R}$: matrices with very few or very many zeros are often factorizable, while those with 3--8 zeros form a ``death valley'' of non-factorizability. We provide rigorous certification using Smale's $\alpha$-criterion and identify specific zero patterns that guarantee factorizability.
\end{abstract}

\section{Introduction}

The Hadamard product $A \circ B$ of two matrices is defined element-wise: $(A \circ B)_{ij} = A_{ij} B_{ij}$. A fundamental question is: given a matrix $C$, when can we find matrices $A$ and $B$ of low rank such that $C = A \circ B$?

This problem arises naturally in:
\begin{itemize}
\item Tensor decomposition and multilinear algebra
\item Algebraic geometry (secant varieties)
\item Quantum information theory
\item Signal processing and compressed sensing
\end{itemize}

By the Schur product theorem, $\rank(A \circ B) \leq \rank(A) \cdot \rank(B)$. Thus, if $C = A \circ B$ with $\rank(A) = \rank(B) = 2$, then $\rank(C) \leq 4$. The converse is our main question: can every rank-4 matrix be expressed as such a product?

\section{The Finite Field Cases}

\subsection{Results over $\mathbb{F}_2$}

We performed an exhaustive search over all $4 \times 4$ matrices over $\mathbb{F}_2$.

\begin{theorem}
Among the 20,160 full-rank $4 \times 4$ matrices over $\mathbb{F}_2$:
\begin{itemize}
\item 14,856 (73.7\%) can be expressed as $A \circ B$ with $\rank(A) = \rank(B) = 2$
\item 5,304 (26.3\%) cannot be so expressed
\end{itemize}
\end{theorem}

The key finding is that matrix density (number of 1s) is 95.7\% predictive of expressibility:

\begin{center}
\begin{tabular}{cc}
\toprule
Number of 1s & Expressible \\
\midrule
16 (zero-free) & 0\% \\
15 & 6.1\% \\
14 & 34.1\% \\
13 & 70.9\% \\
$\leq 12$ & 98--100\% \\
\bottomrule
\end{tabular}
\end{center}

\begin{proposition}[Zero-free obstruction]
No zero-free (all 1s) matrix over $\mathbb{F}_2$ can be expressed as the Hadamard product of two rank-2 matrices.
\end{proposition}

\subsection{Results over $\mathbb{F}_3$}

The $\mathbb{F}_3$ case reveals a dramatic reversal of the $\mathbb{F}_2$ pattern:

\begin{theorem}
Among approximately 811,200 rank-2 $4 \times 4$ matrices over $\mathbb{F}_3$, computational search found 758,012, with the following factorizability by density:

\begin{center}
\begin{tabular}{cc}
\toprule
Number of zeros & Expressibility rate \\
\midrule
0 (zero-free) & 21.0\% \\
1 & 50.0\% \\
2 & 69.2\% \\
3 & 80.8\% \\
4--8 & 84.2\% \\
\bottomrule
\end{tabular}
\end{center}
\end{theorem}

\begin{remark}
The zero-free obstruction from $\mathbb{F}_2$ does not generalize to $\mathbb{F}_3$. In fact, dense matrices become \emph{more} expressible, suggesting the obstruction is specific to characteristic 2.
\end{remark}

\section{The Real Number Case}

\subsection{Algorithmic Approach}

For matrices over $\mathbb{R}$, we cannot use exhaustive search. Instead, we develop an alternating projection algorithm:

\begin{algorithm}[H]
\caption{Alternating Projection for Hadamard Factorization}
\KwIn{Matrix $C \in \mathbb{R}^{4 \times 4}$, tolerance $\varepsilon$}
\KwOut{Matrices $A, B$ of rank $\leq 2$ such that $\|A \circ B - C\|_F < \varepsilon$, or FAILURE}

Initialize $A, B$ with random rank-2 matrices\;
\For{$k = 1$ to MaxIter}{
    $A \leftarrow \arg\min_{\rank(A) \leq 2} \|A \circ B - C\|_F^2$\;
    $B \leftarrow \arg\min_{\rank(B) \leq 2} \|A \circ B - C\|_F^2$\;
    \If{$\|A \circ B - C\|_F < \varepsilon$}{
        \Return $(A, B)$\;
    }
}
\Return FAILURE\;
\end{algorithm}

Each projection step has a closed-form solution using SVD truncation of element-wise quotients.

\subsection{The U-Shaped Factorizability Curve}

\begin{theorem}
For random $4 \times 4$ real matrices with prescribed zero patterns, factorizability follows a U-shaped curve:

\begin{center}
\begin{tabular}{cc}
\toprule
Number of zeros & Factorization success rate \\
\midrule
0 (zero-free) & 48\% \\
1 & 18\% \\
2 & 0\% \\
3--8 & 0--6\% (``death valley'') \\
9 & 3\% \\
10 & 8\% \\
11 & 28\% \\
12 & 50\% \\
\bottomrule
\end{tabular}
\end{center}
\end{theorem}

This U-shape is visualized clearly in UMAP embeddings of the factorizability landscape (see Figure~\ref{fig:umap}).

\subsection{Pattern-Specific Results}

Within each sparsity level, specific zero patterns dramatically affect factorizability:

\begin{theorem}[Pattern Characterization]
The following patterns guarantee high factorizability:
\begin{enumerate}
\item \textbf{Empty rows/columns}: Matrices with $\geq 1$ zero row or column are 15--100\% factorizable
\item \textbf{Block structure}: Matrices with zeros concentrated in off-diagonal blocks show moderate improvement
\item \textbf{Scattered zeros}: Random zero patterns in the 3--8 range are rarely factorizable ($<5\%$)
\end{enumerate}
\end{theorem}

\begin{example}
A matrix with one zero row:
\[
C = \begin{pmatrix}
0 & 0 & 0 & 0 \\
* & * & * & * \\
* & * & * & * \\
* & * & * & *
\end{pmatrix}
\]
is always factorizable as $C = A \circ B$ where the first row of either $A$ or $B$ is zero.
\end{example}

\subsection{Rigorous Certification}

We implement Smale's $\alpha$-criterion for rigorous certification:

\begin{definition}[$\alpha$-criterion]
For the system $F(A,B) = A \circ B - C = 0$, define:
\[
\alpha(A,B) = \beta(A,B) \cdot \|DF(A,B)^{-1} F(A,B)\|
\]
where $\beta(A,B)$ bounds the second derivative. If $\alpha < \alpha_0 = \frac{3 - \sqrt{7}}{2} \approx 0.157671$, then Newton's method converges quadratically to a solution.
\end{definition}

\begin{theorem}[Certification Rate]
Among successfully factorized matrices:
\begin{itemize}
\item 89\% satisfy the $\alpha$-criterion, providing a rigorous certificate
\item 11\% converge numerically but fail certification (typically near rank boundaries)
\end{itemize}
\end{theorem}

\section{Theoretical Implications}

\subsection{Finite Field Phenomena}

The reversal between $\mathbb{F}_2$ and $\mathbb{F}_3$ suggests:

\begin{conjecture}
The zero-free obstruction is specific to characteristic 2. For odd characteristic fields, dense matrices are generically expressible as Hadamard products of low-rank matrices.
\end{conjecture}

\subsection{Real Number Geometry}

The U-shaped curve over $\mathbb{R}$ reveals rich geometric structure:

\begin{proposition}[Dimension Counting]
For a $4 \times 4$ matrix $C$ to equal $A \circ B$ with $\rank(A) = \rank(B) = 2$:
\begin{itemize}
\item Degrees of freedom in $(A,B)$: $2 \times (4 \times 2 + 2 \times 4) - 4 = 28$
\item Constraints from $C$: 16
\item Excess dimensions: 12
\end{itemize}
Thus the variety of expressible matrices has dimension 16 (full), but with special structure near boundaries.
\end{proposition}

\subsection{Connection to Tensor Rank}

Our problem connects to tensor rank decomposition:

\begin{remark}
A matrix $C$ expressible as $A \circ B$ with $\rank(A) = \rank(B) = 2$ corresponds to a tensor in $\mathbb{R}^4 \otimes \mathbb{R}^4$ with CP-rank at most 4.
\end{remark}

\section{Open Questions}

\begin{enumerate}
\item \textbf{Higher dimensions}: Characterize Hadamard factorizability for $n \times n$ matrices with $n > 4$.

\item \textbf{Intermediate fields}: What happens over $\mathbb{F}_p$ for primes $p > 3$?

\item \textbf{Complexity}: Is determining Hadamard factorizability NP-hard in general?

\item \textbf{Algebraic characterization}: Find polynomial invariants that detect non-factorizability.

\item \textbf{Applications}: Exploit the structure for matrix completion or compressed sensing.
\end{enumerate}

\section{Conclusion}

The Hadamard factorization problem reveals surprising richness across different fields. The computational evidence suggests deep connections between matrix sparsity patterns, field characteristics, and geometric properties of secant varieties. The U-shaped factorizability curve over $\mathbb{R}$ and the pattern-specific results provide both theoretical insights and practical algorithms for this fundamental problem in multilinear algebra.

\section*{Acknowledgments}

Computational experiments were performed using JAX for GPU acceleration and rigorous certification via interval arithmetic. Interactive tools available at \url{https://github.com/[repository]}.

\bibliographystyle{plain}
\bibliography{references}

\end{document}